\documentclass[11pt,a4paper]{article}
\usepackage[T1]{fontenc}
\usepackage[utf8]{inputenc}
\usepackage{amsmath,amssymb,xcolor}
\usepackage[margin=25mm]{geometry}
\usepackage[colorlinks=true,urlcolor=blue]{hyperref}
\definecolor{nexorablue}{HTML}{2459B5}
\setlength{\parindent}{0pt}
\setlength{\parskip}{7pt}
\setlength{\emergencystretch}{2em}
\renewcommand{\familydefault}{\sfdefault}
\begin{document}

{\small\bfseries\color{nexorablue} NEXORA CPGE}\par
{\small MPSI / NEX-MES-01 / Corrigé}\par
{\LARGE\bfseries Quatre lectures au banc optique}\par
\vspace{0.25em}\hrule\vspace{0.6em}

\section*{Moyenne et dispersion}
La moyenne est $20{,}15\,\mathrm{cm}$. Les écarts sont $-0{,}01$, $0{,}01$, $0{,}02$, $-0{,}02\,\mathrm{cm}$. La somme de leurs carrés vaut $0{,}0010\,\mathrm{cm^2}$ et $s=\sqrt{0{,}0010/3}\simeq0{,}0183\,\mathrm{cm}$.

\section*{Incertitude de la moyenne}
On divise $s$ par $\sqrt4=2$. Ainsi $u_A\simeq0{,}0091\,\mathrm{cm}$. Avec un chiffre significatif pour $u_A$, le résultat devient $(20{,}150\pm0{,}009)\,\mathrm{cm}$ pour cette contribution seule.


\end{document}
